\section{Logistic speciation}



The number of individuals, or size of a population grows logistically according to standard logistic growth:
$$\ddt{N} = rN\lrb{1-\frac{N}{K}}.$$
The size of population, $N(t)$ is taken to be a continuous variable, a real number. It is assumed that this number of individuals is large compared to the count of species within the total group.\par

Types beget further variations in a descent of speciation that behaves as a Yule pure birth process. That is, given a number of types, $\yt$, initially 1, changes occur as defined: 
$$\pr{\yy_{t+\delt} - \yy_t = 1} = \beta \yy_t \delt + o(\delt)$$
taken in the limit, $\delt\to 0$. In the standard Yule process --- originally specified in 1924 to model the development of species --- $\beta$ is a constant value. Here however, we define this rate of speciation as a function of time. 
$$\beta \equiv \beta(t).$$
Such a step has been made before. Both in the field of evolutionary modeling, and elsewhere. If we were to set $\beta = \ee^{-t}$, the mean of $\yt$ would be described by the deterministic Gompertz model, introduced in 1825 \cite{gompertz1825xxiv}.\par

It doesn't appear that anyone has yet published a version of the model that uses the logistic function. There are a few ways of doing this, for now we will assign $\beta(t)$ the form, 
\begin{equation}
\label{rate1}
\beta(t) = \sigma\lrb{K - N(t)}.
\end{equation}
It serves as an index for the relative abundance of resources experienced by the population.\par

When $N(t)$ is small compared to the carrying capacity of the ecosystem, $K$, $\beta(t)$ is towards the higher end of it's range; species beget species with greather fecundity. We will later explore other options for the birth rate, for example, it may be sensible to ascribe the formula as $\beta(t) = \sigma\cdot \dd N/\dd t$ which would see the highest rate of division at the time of most rapid expansion --- perhaps this makes more sense than the fastest speciation beginng when the population is smallest. But for now let us continue with formula (\ref{rate1}). 

\vspace{10pt}
\hrule
\vspace{10pt}
\noindent
\begin{quotation}
\noindent
\textbf{Aside:} While writing this out, I am starting to wonder if the rate of population expansion is actually orders of magnitude faster than that of speciation. What if the logistic growth doesn't itself spur on type divisions, but instead lays out the potential for them. In this case, it would be incorrect to account for high early rates of speciation with the period of logistic expansion, however, such an expansion may mark the initial condition of a more DDD like model (density dependent diversification) such as the one we were initially working on.  \par

Could it be there is some effect whereby a homogeneos population experiences a saftey in numbers effect, i.e. a larger group has greater control of resorrces and its individuals are more protected from predators and competitors. In this case, the propensity to vary could be inflated. Then as the homogeneos group divides into ever smaller and smaller sub groups, they become more vulnerable --- this leading to harsher selection pressures and lower diversification rate.  
\end{quotation}
\vspace{10pt}
\hrule  
\vspace{10pt}

\noindent
The topic of variable variation in a group has been discussed by many thinkers: from Lamarck at the turn of the 19th centuary, to Darwin in the late 19th centuary; a number of more modern contributers such as S. J Gould in the latter 20th centuary. 
Somewhere in the middle, Ronald Fisher in his 1929 book, The Genetical Theorey of Natural Selection, implies that variations to a type are in constant and abundant supply, and that it is the harshness of environmental conditions in conjunction with sexual mixing which constantly work to quell this significant tendancy for genomes to become irretrievably disprate. 




According to Ronald Fisher, variations to the norm of a given type are in constant and abundant supply. It is the harshness of environmenal conditions which prevents many of them from becoming established; selection pressure supports adhearnce to a current optimal form. But when the selection pressures are lesser, when for example there are more comfortable conditions in the form of a fresh abundance of resources available to the population, variations to the norm are more likely to survive. 
\begin{quotation}
It is known by the experience of breeders that species which recieve superabundant nourishment, and in general a surplus of protection and care, immediately tend in the most marked way to develop variation, and are fertile in prodigies and monstrosities." (Nietzsche beyond good and evil section 262).
\end{quotation}

\begin{quotation}
"To maintain a stationary variance, fresh mutations must be available in each generation." (R. Fisher, the genetical theory of natural selection)
\end{quotation}

The whole thing is reminiscent of the concept described by Ilya Prigogine of dissipative structures. The idea being that where a flow of energy of matter is seen, structures can emerge which apparently go against the second law of thermodynamics. Locally, order can emerge as a spontaneous reduction in entropy. And this emergence is characterised by chaotic dynamics, its exact formation arising influenced by small fluctuations whos effects magnify and propergate.\par
Examples of dissipative structures can be seen throughout the natural world. It seems that one is the evolution of life into a multitude of forms and varieties.
